\section{Background}

\subsection{Finetuning}

Let $x$ represent the input, $h$ represent the hidden state representation, and $W_0$ denote the frozen pretrained weight matrix.

\begin{itemize}
	\item \textbf{Linear probing}: The pretrained model acts as a frozen feature extractor yielding representations $h = f_{W_0}(x)$. Only a linear classification head is trained: \( \hat{y} = W_{\text{head}} h + b \)
	\item \textbf{Full fine-tuning}: All parameters are updated. A full-rank weight update $\Delta W$ is learned: \( h = (W_0 + \Delta W) x \)
	\item \textbf{Parameter-efficient fine-tuning (PEFT)}:
	      \begin{itemize}
		      \item \textbf{LoRA}: The weight update is constrained by decomposing it into two trainable low-rank matrices, $B$ and $A$: \( h = W_0 x + B A x \)
		      \item \textbf{Adapter}: A trainable bottleneck feed-forward network is inserted after the frozen layer's output $h = W_0 x$: \( h' = h + W_{\text{up}} \sigma(W_{\text{down}} h) \)
		      \item \textbf{Prompt tuning}: A trainable sequence of continuous prompt embeddings $P$ is prepended to the input sequence embeddings $X$: \( X' = (P; X) \)
		      \item \textbf{Prefix tuning}: Trainable prefix vectors ($P_K$ and $P_V$) are prepended directly to the Keys ($K$) and Values ($V$) computed by the frozen attention layers: \( h = \text{Attn}(Q, (P_K; K), [P_V; V]) \)
	      \end{itemize}
\end{itemize}

\begin{figure}[htbp]
	\centering
	\includegraphics[width=0.88\textwidth]{training_and_finetuning.png}
	\caption{Training and fine-tuning}
	\label{fig:training_and_finetuning}
\end{figure}

\subsection{Low-Rank Structures}

\subsubsection{Motivation}

Li et al. \cite{li2018measuring} define intrinsic dimension as the codimension of the solution set: if the full parameter space has dimension $D$ and the solution set has dimension $s$, then $D=d_{\mathrm{int}}+s$. Their random-subspace formulation keeps the initialization $\theta_0^{(D)}\in\mathbb R^D$ and projection $P\in\mathbb R^{D\times d}$ fixed while optimizing lower-dimensional coordinates $\theta^{(d)}\in\mathbb R^d$. At the final solution,
\[
	\theta_*^{(D)}=\theta_0^{(D)}+P\theta_*^{(d)}
\]
Figure~\ref{fig:intrinsic_dimension} illustrates the case $D=3$, $s=1$, and $d=d_{\mathrm{int}}=2$. Li et al. establish almost-sure intersection at the codimension threshold for hyperplane solution sets, while arbitrary solution-set topologies need not intersect.

% ==========================================================
% Figure: Intrinsic dimension
% ==========================================================

\begin{figure}[htbp]
	\centering

	\begin{tikzpicture}[
			x=1cm,
			y=1cm,
			>=Stealth,
			font=\small,
			axis/.style={->,semithick},
			update/.style={->,semithick},
			solution/.style={
					line width=1.2pt,
					dash pattern=on .5pt off 3pt,
					line cap=round
				},
			plane/.style={
					fill=gray!25,
					fill opacity=1,
					draw=black!70,
					draw opacity=1,
					line width=.8pt
				},
			point/.style={
					circle,
					fill=black,
					inner sep=1.6pt
				}
		]

		% ------------------------------------------------------
		% (a) Full parameter space
		% ------------------------------------------------------
		\begin{scope}[shift={(0,0)}]

			% Negative theta_1 direction
			\draw[semithick] (-1.25,0) -- (0,0);

			% Positive theta_1 direction
			\draw[axis] (0,0) -- (2.15,0)
			node[right] {$\theta_1$};

			% Negative theta_2 direction
			\draw[semithick] (0,-.5) -- (0,0);

			% Positive theta_2 direction
			\draw[axis] (0,0) -- (0,2.9)
			node[above] {$\theta_2$};

			% Negative theta_3 direction
			\draw[semithick] (1.3,.86) -- (0,0);

			% Positive theta_3 direction
			\draw[axis] (0,0) -- (-1.3,-.86)
			node[below left] {$\theta_3$};

			% Initialization and final solution
			\coordinate (theta0L) at (.78,.65);
			\coordinate (thetaStarL) at (-.58,1.45);

			% Optimization displacement
			\draw[update] (theta0L) -- (thetaStarL);

			% Points
			\node[point] at (theta0L) {};
			\node[point] at (thetaStarL) {};

			% Labels
			\node[below right=2pt] at (theta0L)
			{$\theta_0^{(D)}$};

			\node[above left=2pt] at (thetaStarL)
			{$\theta_*^{(D)}$};

			% Panel heading
			\node[font=\bfseries] at (.35,3.65)
			{(a) Full parameter space};

			\node at (.4,-1.6) {$D=3$};

		\end{scope}


		% ------------------------------------------------------
		% (b) Affine optimization subspace
		% ------------------------------------------------------
		\begin{scope}[shift={(7.2,0)}]

			% --------------------------------------------------
			% Coordinate axes
			% --------------------------------------------------

			% Negative theta_1 direction
			\draw[semithick] (-1.25,0) -- (0,0);

			% Positive theta_1 direction
			\draw[axis] (0,0) -- (2.15,0)
			node[right] {$\theta_1$};

			% Negative theta_2 direction
			\draw[semithick] (0,-.5) -- (0,0);

			% Positive theta_2 direction
			\draw[axis] (0,0) -- (0,2.9)
			node[above right] {$\theta_2$};

			% Negative theta_3 direction
			\draw[semithick] (1.3,.86) -- (0,0);

			% Positive theta_3 direction
			\draw[axis] (0,0) -- (-1.3,-.86)
			node[below left] {$\theta_3$};

			% --------------------------------------------------
			% Opaque two-dimensional affine plane
			% --------------------------------------------------
			\coordinate (A) at (-.95,.05);
			\coordinate (B) at (2.25,.55);
			\coordinate (C) at (.70,2.35);
			\coordinate (E) at (-2.50,1.85);

			\filldraw[plane]
			(A) -- (B) -- (C) -- (E) -- cycle;

			% --------------------------------------------------
			% Initialization and final solution
			% --------------------------------------------------
			\coordinate (theta0R) at (.78,.65);
			\coordinate (thetaStarR) at (-.58,1.45);

			% Entry point of the solution line into the plane
			\coordinate (solutionEntry) at (-1.00687,.11603);

			% --------------------------------------------------
			% One-dimensional solution set
			% --------------------------------------------------

			% Solid portion below the surface
			\draw[semithick]
			(-1.156,-.35) -- (solutionEntry);

			% Sparse dotted portion behind the surface,
			% ending exactly at the final solution
			\draw[solution]
			(solutionEntry) -- (thetaStarR);

			% Solid portion above the final solution
			\draw[semithick]
			(thetaStarR) -- (-.116,2.9);

			% --------------------------------------------------
			% Optimization displacement inside the plane
			% --------------------------------------------------
			\draw[update]
			(theta0R) -- (thetaStarR)
			node[pos=.48,above right,
				fill=white,inner sep=2pt]
			{$P\theta_*^{(d)}$};

			% Points
			\node[point] at (theta0R) {};
			\node[point] at (thetaStarR) {};

			% Point labels
			\node[below right=2pt] at (theta0R)
			{$\theta_0^{(D)}$};

			\node[above left=2pt] at (thetaStarR)
			{$\theta_*^{(D)}$};

			% Panel heading
			\node[font=\bfseries] at (.35,3.65)
			{(b) Affine optimization subspace};
		\end{scope}

	\end{tikzpicture}

	\caption{
		Geometric illustration of intrinsic dimensionality,
		adapted from Li et al.~(2018).
		The one-dimensional solution set $\mathcal S$ is
		represented by a solid line outside the plane and
		a sparse dotted line inside it, intersecting the
		plane at $\theta_*^{(D)}$.
	}
	\label{fig:intrinsic_dimension}

\end{figure}

The low-dimensional fine-tuning results motivate LoRA, but do not establish that weight updates have low matrix rank. Figure~\ref{fig:lora_motivation} summarizes the progression from low-rank observations to LoRA's hypothesis and parameterization.


% ==========================================================
% Figure: Motivation for LoRA
% ==========================================================

\begin{figure}[htbp]
	\centering

	% IMPORTANT: Do not insert blank lines inside this
	% tikzpicture option list.
	\begin{tikzpicture}[
			>=Stealth,
			font=\small,
			node distance=5mm,
			box/.style={
					draw=black!70,
					fill=gray!8,
					rounded corners=2pt,
					text width=12.2cm,
					align=left,
					inner sep=7pt
				},
			hypothesis/.style={
					box,
					dashed,
					fill=gray!3
				},
			method/.style={
					box,
					fill=gray!17
				},
			arrow/.style={
					->,
					semithick
				}
		]

		% --------------------------------------------------
		% 1. Low-rank structure
		% --------------------------------------------------
		\node[box] (n1) {
			\textbf{1. Low-rank structure in machine learning}

			Many learning problems exhibit intrinsic low-rank
			structure. Overparameterized neural networks can
			also develop low-rank properties after training,
			and previous work has explicitly exploited these
			properties through low-rank constraints.

			\hfill{\footnotesize
				Oymak et al. \cite{oymak2019generalization};
				Hu et al. \cite{hu2022lora}
			}
		};


		% --------------------------------------------------
		% 2. Intrinsic dimensionality of fine-tuning
		% --------------------------------------------------
		\node[box,below=of n1] (n2) {
			\textbf{2. Low intrinsic dimensionality of fine-tuning}

			Pretrained language models can be fine-tuned
			effectively through low-dimensional
			reparameterizations of their parameter space.

			\hfill{\footnotesize Aghajanyan et al. \cite{aghajanyan2021intrinsic}
			}
		};


		% --------------------------------------------------
		% 3. LoRA hypothesis
		% --------------------------------------------------
		\node[hypothesis,below=of n2] (n3) {
			\textbf{3. LoRA hypothesis: low-rank weight updates}

			Do downstream tasks require only low-rank
			changes to pretrained weight matrices?

			\[
				\operatorname{rank}(\Delta W)
				\ll \min(m,n).
			\]

			Low-dimensional fine-tuning motivates this
			hypothesis, but does not imply it:
			intrinsic dimension and matrix rank are distinct
			quantities.

			\hfill{\footnotesize Hu et al. \cite{hu2022lora} }
		};


		% --------------------------------------------------
		% 4. LoRA method and consequences
		% --------------------------------------------------
		\node[method,below=of n3] (n4) {
			\textbf{4. LoRA: low-rank parameterization}

			Freeze the pretrained matrix and optimize only
			its low-rank update:
			\[
				W=W_0+\Delta W=W_0+AB,
			\]
			Fewer trainable parameters, reduced optimizer-state
			memory, task-specific adapters, and the ability to
			merge $AB$ into $W_0$ for inference.

			\hfill{\footnotesize Hu et al. \cite{hu2022lora} }
		};


		% --------------------------------------------------
		% Connections
		% --------------------------------------------------
		\draw[arrow] (n1.south) -- (n2.north);

		% Dashed transition: motivation, not implication
		\draw[arrow,dashed] (n2.south) -- (n3.north);

		\draw[arrow] (n3.south) -- (n4.north);

	\end{tikzpicture}

	\caption{
		Conceptual progression motivating Low-Rank Adaptation.
		Observations of low-rank structure in machine learning
		and low intrinsic dimensionality in language-model
		fine-tuning motivate the hypothesis that task-specific
		weight updates have low intrinsic rank.
		The dashed transition indicates a hypothesis rather
		than a mathematical implication.
		LoRA implements this hypothesis through trainable
		low-rank factors while freezing pretrained weights.
	}
	\label{fig:lora_motivation}
\end{figure}


\subsection{LoRA: Low-rank Adaptation}

Hu et al. \cite{hu2022lora} introduce LoRA as a parameter-efficient alternative
to updating all pre-trained weights during downstream fine-tuning.
For a frozen matrix $W_0$, the method learns a low-rank update
$\Delta W=AB$, scaled by $\alpha/r$. The factors have rank dimension
$r$ and are much smaller than the original matrix when
$r\ll\min(d,k)$. This reduces the trainable parameters and the associated
gradient and optimizer state. The frozen model weights still need to be
stored and used for computation. After training, an adapter can be merged
into the base weights, or retained separately when switching among tasks
is useful.

\subsection{O-LoRA: Low-rank Adaptation in Continual Learning}

Continual learning studies how a model can learn a sequence of tasks
without losing useful performance on earlier tasks. Catastrophic
forgetting creates a tension between plasticity, the ability to learn new
tasks, and stability, the retention of prior knowledge. O-LoRA method of
Wang et al. \cite{wang2023olora} assigns a separate LoRA update to each task,
freezes the updates for completed tasks, and adds a regularization term that
encourages the current update subspace to be orthogonal to earlier ones.
Its finite penalty encourages subspace separation; it does not
guarantee exact orthogonality or eliminate forgetting.

\subsection{Relationship and scope}

Both methods use low-rank updates to reduce the part of the model that is
trained. LoRA addresses adaptation to a task, while O-LoRA adds a
continual-learning mechanism intended to reduce interference across a
sequence of tasks. Their reported evaluations use different models,
datasets, and experimental settings, so the papers support a conceptual
comparison of method goals and trade-offs rather than a direct ranking of
performance. Our experiment evaluates LoRA on AG News and does not
evaluate O-LoRA.

\begin{table}[htbp]
	\centering
	\small
	\caption{Conceptual comparison of LoRA and O-LoRA.}
	\label{tab:adapter-method-comparison}
	\begin{tabularx}{\textwidth}{@{}>{\raggedright\arraybackslash}p{2.2cm}XX@{}}
		\toprule
		\textbf{Aspect} & \textbf{LoRA}                                             & \textbf{O-LoRA}                                                                            \\
		\midrule
		Objective       & Adapt a pre-trained model with fewer trainable parameters & Learn sequential tasks while limiting interference efficiently                             \\
		Training        & Train low-rank factors for the selected task              & Train current task factors; freeze earlier task factors                                    \\
		Forgetting      & No explicit forgetting-control mechanism                  & Orthogonality penalty aims to mitigate forgetting                                          \\
		Storage         & One small adapter per retained task                       & Adapter storage grows with the number of tasks                                             \\
		Inference       & Selected adapter may be merged into base weights          & Task updates may be accumulated into weights; factors may be retained for further learning \\
		\bottomrule
	\end{tabularx}
\end{table}
